% ===== §8 =====
\section{The Damage to Generativity}
\label{p25:sec:damage}

\noindent
The second movement of the paper begins here, and its question differs in kind from the first. The first movement asked whether intimate relation furnishes a well-defined optimization problem and answered that it does not; it passed no judgment on optimizing, establishing only that here the operation is undefined. This movement takes that result as settled and asks a question the result makes newly pressing: what happens when the frame is imposed regardless, when a domain that furnishes no well-posed maximization is treated as though it did, its relational and incommensurable and endogenous goods forced into the shape of an objective to be maximized? From here the claims are value-laden, and openly so; and the separation insisted on in the introduction is what earns them the right to be, for they no longer trade on the formal result but build upon it, taking as given that no optimization problem is present and asking after the cost of pretending one is. The first such cost, argued in this section, is a damage with a definite form: the imposition converts a value that should circulate and accumulate within a relation into a quantity that can be settled, extracted, and in the end exhausted.

\subsection*{Two fates for generated value}

The theory of generative justice, developed in this series in exchange with Eglash, distinguishes two fates that the value generated in a relation may meet, and the distinction is the hinge of this section \parencite{eglash2016}. Value may \emph{circulate}: return to those who generate it, replenishing the conditions of its own further generation, so that the relation is a site of ongoing production that sustains itself and the value stays in motion among its makers. Or value may be \emph{extracted}: drawn off from its generators, alienated, settled somewhere outside the circuit of its making, so that the relation is mined rather than cultivated and its capacity to generate is drawn down toward depletion. Eglash's own cases are drawn from labor, ecology, and the technical arts, where the value in question is relatively legible and its diversion relatively easy to see; the contribution of this series has been to carry the distinction into intimate life, where the value at stake is least legible and its diversion hardest to name, and where, this paper now argues, the optimizing frame is one of the chief instruments of the second fate. The claim is not a loose analogy between mining ore and pricing love. It is that the operation optimization performs on a relational good is, in its formal structure, an operation of extraction, and the next step is to show why.

\subsection*{The effect of optimization on a circulation}

To see the extraction one must see what optimization does to a value when it forces it into an objective. To make a relational good a term in $J$ is to render it a \emph{quantity}: a magnitude that can be summed with others, compared, traded off, and, above all, \emph{settled}, brought to a final accounting at which one reads off whether the maximum was reached. But a relational good is not, in its own manner of being, a settleable quantity. It is a circulation, a generation whose life is in its continuing and which has no natural terminus at which it is complete and may be totalled. Meaning between two people is not a stock that accumulates toward a full amount; it is a movement that is either kept up or allowed to lapse. Rendered as a quantity to be maximized, such a good is made available to be settled, and to settle a circulation is to stop it. The optimizing frame does not merely mismeasure the relational good, assigning it a number that is off. It re-forms the good into the kind of thing that can be cashed out, and the cashing out is the depletion, because what is cashed out is no longer in circulation. The harm is thus not that the optimization reaches a wrong figure. It is that the operation of rendering-as-quantity, which optimization must perform before it can begin, has already converted a generation into a stock and thereby positioned it to be drawn down.

\subsection*{A single geometric illustration}

The series has given this depletion a geometric form elsewhere, and it earns one sentence here, strictly as illustration and not as apparatus to be developed. In the account of the good and the bad cycle, generative value was shown to be carried not in any endpoint state but in the accumulated phase of a circuit, the holonomy an interpreter gathers in traversing it, which no terminal reckoning can capture because it is a property of the path and not of any point upon it \parencite{paperix}. The single relevant point is this: that optimization, whose entire grammar is the selection of a best \emph{point}, an $x^{\ast}$ in the feasible set, is by its very form blind to value that lives in a path rather than a point, and so cannot so much as register, let alone protect, the generative accumulation it settles away. The apparatus is not needed here and will not be unfolded; it serves only to name, in a vocabulary the series has already made precise, why a value carried in circulation is invisible to an operation that knows only summits.

\begin{casebox}{the audited relationship}
A couple, anxious and conscientious, begin to keep an account of their relationship: to track what each contributes and receives, to note whether the returns justify the investment, to hold periodic reckonings at which the question is whether, on balance, the thing is paying off. Their intention is care; they mean to tend the relationship by measuring it. But the measuring does to the relationship exactly what this section describes. What had been a circulation, a giving and receiving whose worth lay in its ongoing motion and never stood still to be totalled, becomes a stock to be assessed at each reckoning; and a stock assessed is a stock that can come out short. The damage is not that the audit might return an unfavorable balance and end the relationship, though it might. The damage is prior to any balance and independent of its sign: it is that the auditing has changed the good audited, converting a generation kept alive by not being totalled into a quantity that exists to be totalled, so that even a favorable balance is now a balance struck on a diminished thing. The couple set out to protect the circulation and, by rendering it accountable, stopped it in the act of counting it.
\end{casebox}

\subsection*{The form of the damage}

The damage can now be stated in its proper shape, and stated so that it is plainly a harm and not a mere infelicity of description. When a relational good is subjected to the optimizing frame it undergoes three movements. It is first rendered as a term, which requires assigning it a settleable magnitude it did not possess. It is then entered into an objective to be maximized, which subordinates it to the reckoning of a total in which it figures as one summand among others. And it is then, in the maximizing, spent: treated as a stock to be drawn down toward a maximum rather than a circulation to be kept in motion. This is the formal signature of extraction, and it is why the couple who begin in earnest to ask whether the relationship is worth it have already begun the conversion the theory of generative justice warns against, whatever answer they reach. The harm is not the risk of an unfavorable verdict. It is that the reckoning itself has changed the good reckoned, in the same way and for the same reason that the pricing of the previous part diminished the friendship it priced: by treating as a settleable quantity what had its life in never being finally settled. Optimization imposed where no optimization problem exists is thus not an idle error that leaves its object untouched while failing to improve it. It is an operation that remakes its object into something it can act upon, a settleable quantity, and in the remaking performs the very extraction that draws off into a final accounting the value whose life was in never being finally accounted.
